\documentclass[10pt,journal,compsoc]{IEEEtran}
% IEEE Computer Society の学術誌（TAC）向けの IEEEtran の設定（bare_jrnl_compsoc に相当）．
% 雛形は最終の体裁を表さない（公式の注意）ので，投稿時の PDF は 12 頁より短く（11.5 頁を目安）に保つ．
\usepackage[T1]{fontenc}
\usepackage[utf8]{inputenc}
\usepackage{amsmath,amssymb}
\usepackage{graphicx}
\usepackage{booktabs}
\usepackage{multirow}
\usepackage[nocompress]{cite}
\usepackage{url}
\usepackage[hidelinks]{hyperref}
\newtheorem{proposition}{Proposition}

% 本文の実験由来の数値は \NUM{key} で書く（results/numbers.json → numbers.tex）．
% 自動生成．手で編集しない．
% 生成元: results/numbers.json ／ 生成器: scripts/gen_numbers_tex.py
% 本文の実験由来の数値は \NUM{key} で書く．未登録の key はビルド時にエラーになる．
\makeatletter
\providecommand{\NUM}[1]{%
  \@ifundefined{NUM@#1}{%
    \GenericError{}{numbers.tex: key '#1' is not registered in results/numbers.json}{}{}%
    \textbf{??}%
  }{\csname NUM@#1\endcsname}%
}
\makeatother


% 図：figures/<name>.pdf が無いうちは枠だけを出す（図の台本は scripts/figures/fig_<name>.py）．
\newcommand{\FIG}[2]{\IfFileExists{figures/#1.pdf}{\includegraphics[width=#2]{figures/#1.pdf}}{\fbox{\parbox[c][6em][c]{\dimexpr#2-2\fboxsep-2\fboxrule\relax}{\centering\footnotesize figure \texttt{#1} (pending)}}}}


\begin{document}

\title{What Few Target Labels Fix First: Threshold-versus-Direction Analysis and Cross-Modal Weak-Label Transfer for Cross-Corpus Multimodal Emotion Recognition}

\author{Author~Names% TODO
\IEEEcompsocitemizethanks{\IEEEcompsocthanksitem Affiliations. TODO}}

\markboth{IEEE Transactions on Affective Computing}%
{Author \MakeLowercase{\textit{et al.}}: What Few Target Labels Fix First}

\IEEEtitleabstractindextext{%
\begin{abstract}
A recognizer moved to a new corpus often has a few labels per emotion plus labeled data from other corpora. We analyze what these labels repair in a linear classifier on frozen audio, face, and text features across all ordered pairs of seven corpus groups. Splitting each pairwise decision into a threshold and a direction, we find that for single-modality audio and face with at most five labels per emotion, a fraction of \NUM{p1-e29-critshare-kle5-single-lo} to \NUM{p1-e29-critshare-kle5-single-hi} of what the labels add to a pooled classifier comes from correcting thresholds, whereas audio directions improve at a rate set by an effective dimension of \NUM{p1-qdir-audio-lo} to \NUM{p1-qdir-audio-hi}, and one source corpus is worth \NUM{p1-e32-n1src-audio-lo} to \NUM{p1-e32-n1src-audio-hi} labels per emotion of direction. Neither few labels nor label-free selection repairs the direction of a strong start. We therefore train an anchored linear student on weak labels that a fused classifier trained on the sources alone, seeing modalities the student lacks, assigns to all unlabeled target utterances. Extended with teachers refit on the target pool, the method beats the strongest of nine identically tuned competitors for all eleven students, including second encoders, at zero and five labels by \NUM{cmp-oursx-vs-best-range-lo} to \NUM{cmp-oursx-vs-best-range-hi} UAR points.
\end{abstract}

\begin{IEEEkeywords}
Multimodal emotion recognition, cross-corpus, few-shot learning, domain adaptation, knowledge distillation, weak labels.
\end{IEEEkeywords}}

\maketitle
\IEEEdisplaynontitleabstractindextext
\IEEEpeerreviewmaketitle

\section{Introduction}
\label{sec:introduction}
\IEEEPARstart{T}{his} is a placeholder. TODO: three research questions, then four numbered contributions (see CLAIM.md~6.2 in docs/paper-claim.md); one number in the contribution sentence.

\section{Related Work}
\label{sec:related}

\emph{Cross-corpus emotion recognition with few target labels.}
Cross-corpus evaluation has long shown large drops in speech emotion recognition and has motivated corpus normalization~\cite{schuller2010crosscorpus}.
Methods that use a few labeled target utterances include siamese few-shot learning for spontaneous speech~\cite{feng2023fewshot}, adversarial domain generalization evaluated with growing amounts of target labels~\cite{gideon2021addog}, pre-finetuning on source corpora~\cite{chen2023prefinetuning}, adaptation between acted and improvised speech and video~\cite{li2021acted}, and active learning on frozen speech features with a logistic regression~\cite{xu2024crosstask}.
Speech-and-text models have been adapted with five to sixty labels per class~\cite{he2025m4ser} or a few labels per corpus across languages~\cite{zaidi2024enhancing}.
These studies propose or evaluate methods; none measures which part of the classifier the target labels change, and none varies $K$ over many corpus pairs with three modalities including the face; recent work on cross-corpus speech-and-text recognition also notes that cross-corpus multimodal methods are few~\cite{ahn2025multitask}.
Gideon et al.~\cite{gideon2021addog} list a formalization of when to generalize and when to specialize as future work, and Feng and Chaspari~\cite{feng2023fewshot} list the effect of source properties and of combining sources; the threshold--direction analysis and the label-equivalent value of a source address these questions for linear heads.
For test-time adaptation of speech emotion recognition without labels, entropy minimization and pseudo-labeling were found to fail in most settings~\cite{dong2026test}, which agrees with our comparison.

\emph{Error of a linear rule estimated from few samples.}
The asymptotic expansion of the misclassification rate of the sample linear discriminant contains terms of order $1/N$ that do not depend on the dimension and terms proportional to it~\cite{okamoto1963asymptotic}, and Raudys expressed the second in terms of an effective dimension~\cite{raudys1998intrinsic}; our threshold cost $1/(2K)$ and direction cost $q/K$ follow this split for a penalized logistic regression, with $q$ measured.
The accuracy formula \eqref{eq:twocloud} is standard for a linear rule on two Gaussian classes~\cite{desilva2023value,dobriban2018highdimensional}.
For frozen features with a linear head, domain-adjusted regression bounds the excess risk through the error of the target mean and a residual~\cite{rosenfeld2022domainadjusted}, linear probing is preferred over fine-tuning under distribution shift~\cite{kumar2022lpft}, and additive logit corrections lose their effect in the limit of vanishing regularization~\cite{kini2021labelimbalanced}.
Finite-sample analyses of transfer with few target samples exist for regression~\cite{ha2026labeled} and for one-dimensional discriminants~\cite{desilva2023value}; we measure the corresponding quantities for multiclass logistic regression on real corpus pairs.

\emph{Value of a source in target samples.}
The value of pretraining or of a source has been expressed as an effective amount of target data in characters~\cite{hernandez2021scaling}, in training examples~\cite{lescao2021data}, as an effective sample size of a prior~\cite{wiesenfarth2020quantification}, and as a rate exponent~\cite{hanneke2019value}.
Our expression \eqref{eq:value} inverts the classical relation between the cosine of an estimated direction and the sample size~\cite{raudys1998intrinsic,dobriban2018highdimensional} to express a source direction as labels per emotion, and we measure it on corpus pairs.

\emph{Pooling and multiple sources.}
Bounds for learning from several domains let the distance between domains act only as a penalty~\cite{bendavid2010theory,blitzer2007learning}.
The bias of pooled linear classifiers has been solved in high-dimensional limits for centered multi-task models with task-specific intercepts~\cite{tiomoko2021deciphering} and for pooled least squares without intercepts~\cite{yang2025precise}, and adding domain information explicitly can beat pooling~\cite{zhu2026domain}.
One Newton step approximates the influence of a group of points~\cite{pregibon1981logistic,koh2019accuracy}.
Applied to added target points without a target intercept, it yields \eqref{eq:lambda}, which extends the one-dimensional threshold shift of De Silva et al.~\cite{desilva2023value} with a weight path through which the separability of the target from the source corrects the threshold.

\emph{Privileged information and cross-modal distillation.}
When tasks share labels and differ in input modality, the extra modality is privileged information~\cite{vapnik2009new,lopezpaz2016unifying}.
Cross-modal distillation transfers supervision between paired modalities~\cite{gupta2016cross}; in emotion recognition, face teachers train speech students on unlabeled video~\cite{albanie2018emotion}, fused speech-and-text teachers train speech students~\cite{srinivasan2022representation}, and audio-visual or video teachers pass soft labels or embeddings to speech students~\cite{pendyala2025leveraging,hajavi2024audio}.
Unimodal teachers can improve fused students~\cite{du2021improving,xue2021multimodal,rakib2025g2d}, and fused outputs can train unimodal heads~\cite{jaritz2020xmuda}.
In these studies the teacher is trained on the target corpus or on a third corpus, or the student uses all modalities at inference.
Our basic teacher is the fused linear classifier of the configuration trained on the source corpora alone, which labels the target pool without being refit; the student is a linear head anchored to a source start and uses only its own modalities at test time. The extended version adds candidate teachers trained on the target pool.

\emph{Weak labels, class balance, anchors, and selection.}
Fixed teachers have labeled all unlabeled data with soft labels for representation learning across tasks with different classes~\cite{phoo2021selftraining} and for transductive few-shot classification with a fixed weight~\cite{zanella2024boosting}, and modality-specific teachers have been distilled into a single modality across domains~\cite{hatano2024multimodal}.
Balancing pseudo-labels across classes is common~\cite{huang2022unsupervised,xie2020selftraining,hu2021leveraging,berthelot2022adamatch}, but class imbalance in the unlabeled set can make such balancing harmful~\cite{veilleux2021realistic}.
Linear probes anchored to a zero-shot classifier with a quadratic penalty~\cite{silvarodriguez2024closer} and hypothesis transfer~\cite{kuzborskij2017fast,li2018explicit} anchor the weights to a prior solution, and leave-one-domain-out validation selects hyperparameters without target labels~\cite{gulrajani2021domainbed}.
Our method combines these components; what is specific to it is a fused teacher trained on sources alone that sees modalities the student lacks, weak labels used to train the classifier itself because sources and target share the emotion set, separate anchors for weights and intercepts, and the selection of the teacher and the weak-label weight for each $K$.

\emph{Theory of weak-label gains and of label-free selection.}
The gain of a student trained on a weak teacher's labels has been bounded from below by the misfit between them~\cite{charikar2024quantifying} and related to the disagreement between their representations~\cite{dong2025discrepancies,xue2025representations}; self-training is analyzed under gradual shift~\cite{kumar2020understanding}.
In Gaussian linear models the optimal design does not depend on the observed responses~\cite{chaloner1984optimal,chaloner1995bayesian}, and with small budgets random or typical selection is hard to beat~\cite{hacohen2022active}.
Proposition~\ref{prop:span} follows from the representer theorem~\cite{kimeldorf1971results,scholkopf2001generalized}.
We use these results for a student anchored to a start: the teacher's UAR advantage over the start sets the scale of the gain, and, under a Gaussian approximation and without a preferred direction of the error, label-free selection cannot align the corrected subspace with the start's direction error.

\section{Setting and Measurement Protocol}
\label{sec:setting}
TODO: seven corpora, four emotions, speaker-disjoint splits; frozen encoders and a linear head; the three training recipes of the map; definition of $K$ (per class); the two-cloud model and $d$, $o$; pseudo-target selection; why $C$ is fixed to the source-only value; practical versus oracle values.

\section{Part I: What Few Target Labels Fix First}
\label{sec:part1-what-few-labels-fix}

Part~I asks two questions about the linear head of Section~\ref{sec:setting}: what $K$ labels per emotion change in it, and how much of that change a source corpus can supply instead.
We first describe the gain map, then state the assumptions under which a threshold and a direction have different estimation costs, and then test four predictions that follow from these costs.

\subsection{The Gain Map}
\label{sec:part1-map}

Adding a source helps when target labels are scarce and hurts when they are plentiful (Fig.~\ref{fig:gainmap}).
For audio, the mean gain at $K=1$ is $\NUM{map-gain-a-k1}$ points, and its 95\% interval lies above zero for \NUM{map-gain-pos-a-k1} of the \NUM{map-pairs-a} ordered pairs.
The mean gain falls below zero by $K=25$ ($\NUM{map-gain-a-k25}$ points), and with the full pool the interval lies below zero for \NUM{map-gain-neg-a-kall} pairs.
All eight configurations follow the same course: the mean gain at $K=1$ ranges from $\NUM{map-gain-v-k1}$ (V) to $\NUM{map-gain-avt-k1}$ points (AVT), and at $K=25$ it is negative in every configuration.
The size of the gain is largely a property of the two corpora.
An additive model with one effect per target and one per source explains \NUM{map-share-a-lo}--\NUM{map-share-v-hi}\% of the variance of the gains and of the largest budget at which a source still helps, across A, A~(HuBERT), V, and AV.
The rest of Part~I explains this course through what the labels change in the classifier.

\begin{figure*}[t]
\centering
\FIG{gainmap}{\textwidth}
\caption{The gain map. (a)~Mean gain of pooling over training on the target labels alone (pooled minus target-only UAR in the same iteration), averaged over ordered pairs, against the label budget $K$ per emotion. (b)~Share of pairs whose 95\% interval of the gain lies above zero, includes zero, or lies below zero; these intervals treat the thirty iterations as independent, which makes them too narrow for IEMOCAP targets. All ordered pairs, including the two pairs of IEMOCAP sessions; the pooled classifier uses the penalty of the source-only classifier.}
\label{fig:gainmap}
\end{figure*}

\subsection{Two Costs of a Few Labels}
\label{sec:part1-costs}

We analyze the pairwise decisions of Section~\ref{sec:setting-model} under three assumptions.
\begin{enumerate}
\item[(A1)] Within each emotion, the score difference $z$ of a pair is Gaussian with a variance shared by the two emotions, so that \eqref{eq:twocloud} gives the pairwise accuracy.
\item[(A2)] The learner keeps a finite $\ell_2$ penalty on the weights and leaves the intercepts unpenalized. With a vanishing penalty on separable data, additive corrections of the logits lose their effect~\cite{CITE:kini2021}; a finite penalty keeps the intercepts free to move.
\item[(A3)] The $K$ labels per emotion are a random draw from the pool.
\end{enumerate}

\emph{Threshold.}
With the direction fixed, the intercept that balances two emotions places the threshold at the midpoint of the two class means of $z$.
Estimated from $K$ labels per emotion, this midpoint has variance $\sigma^2/(2K)$, so the residual offset satisfies $\mathbb{E}[o^2] = 1/(2K)$.
Expanding \eqref{eq:twocloud} to second order in $o$ gives the expected threshold loss that remains,
\begin{equation}
\label{eq:thrcost}
\mathbb{E}\left[a(d, 0) - a(d, o)\right] \approx \frac{d}{4}\,\phi\!\left(\frac{d}{2}\right)\frac{1}{2K} \le \frac{\phi(1)}{4K},
\end{equation}
where $\phi$ is the standard normal density and the maximum over $d$ is attained at $d=2$.
A handful of labels therefore removes most of any threshold offset, however large the offset was.

\emph{Direction.}
Estimating the direction requires estimating a vector, and its error shrinks with the ratio of an effective dimension $q$ to $K$.
For a linear rule estimated from $K$ labels per emotion, the separation behaves as
\begin{equation}
\label{eq:dircost}
d(K) = \frac{\Delta}{\sqrt{1 + q/K}},
\end{equation}
the form that the asymptotic expansion of the misclassification rate of the linear discriminant takes in its direction term~\cite{CITE:okamoto1963,CITE:raudys1998}.
We measure $q$ by fitting \eqref{eq:dircost} to the separations of target-only classifiers along the label budget.
The fitted $q$ is $\NUM{p1-qdir-a}$ for WavLM audio, $\NUM{p1-qdir-ah}$ for HuBERT audio, $\NUM{p1-qdir-av}$ for AV, and $\NUM{p1-qdir-v}$ for the face, so for audio a direction comparable to the target's best needs on the order of twenty labels per emotion.
In the ranges of $K$ studied here, labels thus fix thresholds cheaply and audio directions slowly.
Four predictions follow, and Table~\ref{tab:part1} collects the measurements that test them.
(P1)~Most of the loss of a source-only classifier that a few labels recover is threshold loss.
(P2)~The share of the gain that comes from the threshold falls as $K$ grows.
(P3)~Pooling corrects the threshold in proportion to the weight that the target receives (Section~\ref{sec:part1-pooling}).
(P4)~A source substitutes for a number of target labels set by $q$ and by how well its direction serves the target (Section~\ref{sec:part1-value}).

\begin{figure}[t]
\centering
\FIG{twocloud}{\columnwidth}
\caption{Threshold and direction of a pairwise decision. (a)~Two Gaussian classes along the score difference $z$; the threshold offset $o$ moves the decision point away from the midpoint, and the separation $d$ is set by the direction. (b)~Loss of source-only classifiers against the target's best linear rule, split into the threshold loss and the direction loss, in points of pairwise accuracy, averaged over the six emotion pairs and over ordered pairs.}
\label{fig:twocloud}
\end{figure}

\subsection{What the Labels Fix}
\label{sec:part1-fix}

The two-cloud model reproduces the pairwise accuracies of the classifiers we train.
Over all emotion pairs and ordered pairs, the accuracy predicted from the measured $d$ and $o$ correlates with the measured accuracy at \NUM{p1-e13-fit-r-lo}--\NUM{p1-e13-fit-r-hi}, with a mean absolute error of \NUM{p1-e13-fit-mae-lo}--\NUM{p1-e13-fit-mae-hi} points.

\emph{Source-only classifiers lose mainly through the direction, except on the face.}
For WavLM audio, the threshold loss is \NUM{p1-e13-loss-offset-a} points and the direction loss \NUM{p1-e13-loss-dir-a} points (Fig.~\ref{fig:twocloud}b); for the face the two are \NUM{p1-e13-loss-offset-v} and \NUM{p1-e13-loss-dir-v} points, and for text \NUM{p1-e13-loss-offset-t} and \NUM{p1-e13-loss-dir-t} points.
The threshold loss accounts for only \NUM{p1-e29-lossshare-lo}--\NUM{p1-e29-lossshare-hi} of the total loss across configurations.

\emph{The threshold loss is mostly a shift of the corpus mean.}
Centering, which subtracts the mean of the unlabeled utterances of each corpus and uses no label, raises the source-only UAR by \NUM{p1-center-srconly-a} points for WavLM audio, \NUM{p1-center-srconly-ah} for HuBERT audio, and \NUM{p1-center-srconly-v} for the face, and it raises it for \NUM{p1-center-srconly-up-a-n} of the \NUM{p1-center-srconly-pairs-a-n} audio pairs (three iterations).
Refitting only the intercepts of the source-only classifier on $K=5$ labels per emotion, with the weights frozen, gains \NUM{p1-d1-gain-a-k5} points for WavLM audio, close to the \NUM{p1-merged-gain-a-k5} points that pooling gains with the same labels.
After centering, the two gains shrink to \NUM{p1-center-d1-gain-a-k5} and \NUM{p1-center-merged-gain-a-k5} points (Fig.~S2 of the supplementary material).

\emph{The labels go to the threshold first (P1, P2).}
Writing the improvement of the pooled classifier over the source-only classifier as the sum of a threshold part and a direction part, computed from the measured $d$ and $o$ with \eqref{eq:twocloud}, the threshold part dominates for single-modality audio and face at $K \le 5$: its share is \NUM{p1-e29-critshare-kle5-single-lo}--\NUM{p1-e29-critshare-kle5-single-hi} across WavLM audio, HuBERT audio, and the face at $K=1$ and $K=5$ (Fig.~\ref{fig:thresholdshare}a and Table~\ref{tab:part1}; the shares exclude the two pairs of IEMOCAP sessions).
The share is much larger than the threshold's share of the loss, so the labels correct the cheap part of the error first.
As predicted, the share falls with $K$; at $K=10$ it is \NUM{p1-e29-critshare-k10-ah} for HuBERT audio.
In the fused and text configurations the threshold loss is smaller, and at $K=5$ the threshold part is about half of the gain or less (\NUM{p1-e29-critshare-k5-fused-text-lo}--\NUM{p1-e29-critshare-k5-fused-text-hi}, with \NUM{p1-e29-critshare-k5-av} for AV); at $K=1$ it still exceeds half for AV (\NUM{p1-e29-critshare-k1-av}) and VT (\NUM{p1-e29-critshare-k1-vt}).
A two-layer perceptron on the same features gives threshold shares of \NUM{p1-e30-critshare-k5-audio-lo}--\NUM{p1-e30-critshare-k5-audio-hi} for audio and \NUM{p1-e30-critshare-k5-v-wd2}--\NUM{p1-e30-critshare-k5-v-wd4} for the face at $K=5$, so the pattern is not specific to the linear head, although for HuBERT audio it falls below half.

\emph{The direction stays where the source left it.}
Under \eqref{eq:dircost}, even the optimal combination of the source direction with ten labels per emotion recovers \NUM{p1-dirloss-recov-bound-k10-audio-lo}--\NUM{p1-dirloss-recov-bound-k10-audio-hi}\% of the direction loss for audio, and hypothesis transfer, which starts from the source classifier and refits it on the labels, recovers \NUM{p1-dirloss-recov-ht-k10-audio-lo}--\NUM{p1-dirloss-recov-ht-k10-audio-hi}\%.

\begin{figure}[t]
\centering
\FIG{thresholdshare}{\columnwidth}
\caption{Where the labels go. (a)~Share of the improvement of the pooled classifier over the source-only classifier that comes from the threshold, against $K$ (median over ordered pairs without the two pairs of IEMOCAP sessions; ten iterations). Solid lines: single-modality audio and face; dashed lines: fused and text configurations; the horizontal black line marks one half. (b)~Expected residual errors after estimation from $K$ labels per emotion: the threshold offset falls as $1/(2K)$ and the direction error as $q/K$, with $q$ fitted per configuration.}
\label{fig:thresholdshare}
\end{figure}

\begin{table}[t]
\caption{Part~I quantities per configuration. Losses of the source-only classifier in points of pairwise accuracy; $q$ is the effective dimension of the direction in \eqref{eq:dircost}; the last three columns give the share of the pooled improvement that comes from the threshold (median over ordered pairs).}
\label{tab:part1}
\centering
\footnotesize
\setlength{\tabcolsep}{3.5pt}
\begin{tabular}{@{}lrrrrrr@{}}
\toprule
 & \multicolumn{2}{c}{Loss (points)} & & \multicolumn{3}{c}{Threshold share} \\
\cmidrule(lr){2-3}\cmidrule(l){5-7}
Config. & Thr. & Dir. & $q$ & $K{=}1$ & $K{=}5$ & $K{=}10$ \\
\midrule
A & \NUM{p1-e13-loss-offset-a} & \NUM{p1-e13-loss-dir-a} & \NUM{p1-qdir-a} & \NUM{p1-e29-critshare-k1-a} & \NUM{p1-e29-critshare-k5-a} & \NUM{p1-e29-critshare-k10-a} \\
A (HuBERT) & \NUM{p1-e13-loss-offset-ah} & \NUM{p1-e13-loss-dir-ah} & \NUM{p1-qdir-ah} & \NUM{p1-e29-critshare-k1-ah} & \NUM{p1-e29-critshare-k5-ah} & \NUM{p1-e29-critshare-k10-ah} \\
V & \NUM{p1-e13-loss-offset-v} & \NUM{p1-e13-loss-dir-v} & \NUM{p1-qdir-v} & \NUM{p1-e29-critshare-k1-v} & \NUM{p1-e29-critshare-k5-v} & \NUM{p1-e29-critshare-k10-v} \\
AV & \NUM{p1-e13-loss-offset-av} & \NUM{p1-e13-loss-dir-av} & \NUM{p1-qdir-av} & \NUM{p1-e29-critshare-k1-av} & \NUM{p1-e29-critshare-k5-av} & \NUM{p1-e29-critshare-k10-av} \\
T & \NUM{p1-e13-loss-offset-t} & \NUM{p1-e13-loss-dir-t} & -- & \NUM{p1-e29-critshare-k1-t} & \NUM{p1-e29-critshare-k5-t} & \NUM{p1-e29-critshare-k10-t} \\
VT & \NUM{p1-e13-loss-offset-vt} & \NUM{p1-e13-loss-dir-vt} & -- & \NUM{p1-e29-critshare-k1-vt} & \NUM{p1-e29-critshare-k5-vt} & \NUM{p1-e29-critshare-k10-vt} \\
AT & \NUM{p1-e13-loss-offset-at} & \NUM{p1-e13-loss-dir-at} & -- & \NUM{p1-e29-critshare-k1-at} & \NUM{p1-e29-critshare-k5-at} & \NUM{p1-e29-critshare-k10-at} \\
AVT & \NUM{p1-e13-loss-offset-avt} & \NUM{p1-e13-loss-dir-avt} & -- & \NUM{p1-e29-critshare-k1-avt} & \NUM{p1-e29-critshare-k5-avt} & \NUM{p1-e29-critshare-k10-avt} \\
\bottomrule
\end{tabular}
\end{table}

\subsection{How Pooling Corrects the Threshold}
\label{sec:part1-pooling}

Pooling corrects the threshold without a separate intercept for the target, and one Newton step from the source-only solution shows how (P3).
Let $N_s$ and $N_t$ be the numbers of source and target training utterances, $\kappa_s$ and $\kappa_t$ the mean curvatures of the loss at the source-only solution, $A_s$ the curvature-weighted scatter of the source features, and $\boldsymbol{\delta}$ the difference between the target mean and the curvature-weighted source mean.
One step leaves the fraction $1 - \lambda$ of the source-only offset $o_s$, with
\begin{equation}
\label{eq:lambda}
\lambda = \frac{\rho}{1 + \rho},\qquad \rho = N_t\kappa_t\left[\frac{1}{N_s\kappa_s} + \boldsymbol{\delta}^{\top}\left(A_s + \frac{I}{C}\right)^{-1}\boldsymbol{\delta}\right].
\end{equation}
The first term in the bracket is the intercept path, through which the target corrects the threshold in proportion to its share of the training weight.
The second term is a weight path: when the target is linearly separable from the source along directions in which the source varies little, the pooled classifier places weight on such a direction, and that weight acts as an intercept that only the target utterances see.
This form is the group influence of a set of points~\cite{CITE:koh2019,CITE:pregibon1981} written for added points, and its intercept path reduces to the threshold shift $m/(n+m)$ of a one-dimensional discriminant pooled from $n$ target and $m$ source points~\cite{CITE:desilva2023}.

The measured correction $\lambda = 1 - o_m/o_s$ of pooled classifiers, computed from their $d$ and $o$ on the target, is ordered as \eqref{eq:lambda} predicts, although the expression overestimates its average (Fig.~S1 of the supplementary material).
Over ordered pairs and $K \in \{1, 5, 10\}$, the rank correlation between the predicted and the measured $\lambda$ is \NUM{p1-e15-lam-rho-lo}--\NUM{p1-e15-lam-rho-hi} across A, A~(HuBERT), V, and AV (ten iterations).
The full Newton step orders the corrections with a rank correlation of \NUM{p1-e16-lam-rho-lo}--\NUM{p1-e16-lam-rho-hi}, and it also predicts the change of the direction for V and AV (\NUM{p1-e16-dir-rho-video-lo}--\NUM{p1-e16-dir-rho-video-hi}), but less well for audio (\NUM{p1-e16-dir-rho-audio-lo}--\NUM{p1-e16-dir-rho-audio-hi}), where the second-order expansion is less accurate.
The weight path explains why the distance between corpora does not predict the gain: a large distance lowers the source-only accuracy but also widens the path through which pooling corrects the threshold.
The derivation and the simulations that check \eqref{eq:lambda} are in the supplementary material.

\subsection{What a Source Is Worth}
\label{sec:part1-value}

A source contributes a direction, and its value can be expressed in target labels (P4).
Let $e$ be the efficiency of the source direction for the target, the factor by which it scales the target's best separation, $d = e\Delta$.
Under \eqref{eq:dircost}, the direction estimated from $K$ labels has a noise-to-signal ratio $q/K$, and the source direction has the ratio $(1 - e^2)/e^2$.
Combining two independent estimates with optimal weights adds their sample sizes, so starting from the source and refitting on $K$ labels behaves like training on the target alone with $K + n$ labels per emotion, where
\begin{equation}
\label{eq:value}
n = q\,\frac{e^2}{1 - e^2}.
\end{equation}
Equation~\eqref{eq:value} inverts the classical relation between the cosine of an estimated direction and the sample size~\cite{CITE:raudys1998,CITE:dobriban2018}; the inversion turns the source's direction into a count of labels, and $n$ does not depend on $K$.

We measure $n$ by locating the UAR of hypothesis transfer on the target-only learning curve and subtracting $K$ (Fig.~\ref{fig:sourcevalue}).
For a single audio source, $n$ is \NUM{p1-e32-n1src-audio-lo}--\NUM{p1-e32-n1src-audio-hi} labels per emotion and stays nearly constant from $K=5$ to $K=25$; for the face it is \NUM{p1-e32-n1src-v-k5}--\NUM{p1-e32-n1src-v-k10}.
Starting from all sources pooled, $n$ is \NUM{p1-e32-nall-a-k5} and \NUM{p1-e32-nall-a-k10} labels per emotion for WavLM audio at $K=5$ and $K=10$.
Across audio pairs, the rank correlation between the predicted and the measured $n$ is \NUM{p1-e32-rho-audio-lo}--\NUM{p1-e32-rho-audio-hi}, and the measured values are \NUM{p1-e32-ratio-audio-lo}--\NUM{p1-e32-ratio-audio-hi} times the predicted ones; for all sources pooled the ratio is \NUM{p1-e32-allsrc-ratio-lo}--\NUM{p1-e32-allsrc-ratio-hi}.

\begin{figure}[t]
\centering
\FIG{sourcevalue}{\columnwidth}
\caption{What a source is worth. Predicted value $n = q e^2/(1-e^2)$ against the value measured on the target-only learning curve, in labels per emotion, for single sources and for all sources pooled ($K=5$ and $K=10$; pairs without the two pairs of IEMOCAP sessions). The diagonal marks agreement; points beyond the range of the learning curve are shown in the top band and points with $n \le 0$ in the bottom band.}
\label{fig:sourcevalue}
\end{figure}

\subsection{Summary of Part~I}
\label{sec:part1-summary}

For single-modality audio and face features, a few target labels fix the thresholds, and centering removes most of the threshold error without any label.
The direction is the expensive part: audio needs on the order of $q$ labels per emotion to approach the target's best direction, and a source supplies about ten of them, all sources together two to three times as many.
Further progress with $K \le 5$ therefore requires more direction information than the labels and the sources provide in this form.
Part~II derives where that information can come from.

\section{Part II: From the Mechanism to a Method}
\label{sec:part2-mechanism-to-method}

Part~I leaves the following situation: centering and a refit of the intercepts handle the threshold, a source supplies about ten labels' worth of direction, and the direction error that remains is the main loss.
Part~II builds the strongest start that sources allow, shows that neither a few labels nor a label-free choice of which utterances to label can move its direction much, and then supplies the missing direction information through weak labels on the whole unlabeled pool, from a teacher that sees another modality (Fig.~\ref{fig:method}).

\begin{figure*}[t]
\centering
\FIG{method}{\textwidth}
\caption{The method. Source and target features are centered separately. All sources outside the target corpus form a start classifier $(W_0, b_0)$, and the student starts from it with separate $\ell_2$ anchors on the weights and the intercepts. A teacher trained on the sources alone, for a single-modality student typically the fused classifier of the other modalities, labels every unlabeled target utterance with soft labels whose per-emotion totals are equalized; these weak labels enter the student's loss with total weight $m$ next to the $4K$ true labels. The teacher, $m$, and the anchors are selected for each $K$ on pseudo-targets.}
\label{fig:method}
\end{figure*}

\subsection{The Start: Centering and Anchored Transfer}
\label{sec:part2-start}

The student operates on centered features: each source is centered with its own mean, and the target with the mean of its unlabeled pool, which removes most of the threshold loss (Section~\ref{sec:part1-fix}).
The start $(W_0, b_0)$ pools the directions of all sources outside the target corpus.
Each source $s$ contributes the closed-form discriminant direction $W_s(\rho_s) = ((1-\rho_s)\Sigma_s + \rho_s \tau_s I)^{-1}\tilde{M}_s$, where $\tilde{M}_s$ holds its centered class means, $\Sigma_s$ is its covariance, and $\tau_s$ the mean of its eigenvalues; the directions are summed with weights proportional to $N_s^{\gamma}$, and the intercepts assume equal class priors.
The student solves
\begin{equation}
\label{eq:student}
\min_{W, b}\ \sum_{i \in \mathcal{L}} \ell\!\left(W \mathbf{z}_i + b, y_i\right) + \frac{\|W - W_0\|_{M}^2}{2C'} + \frac{\|b - b_0\|^2}{2C_b},
\end{equation}
where $\mathcal{L}$ holds the $4K$ labeled utterances and the norm $\|\cdot\|_M$ uses a metric $M$ built from the scatter of the unlabeled target pool with mixing weight $\rho_M$.
Separate anchors for the weights and the intercepts follow Part~I: the intercepts need little anchoring because a few labels set them well, while the weights need strong anchoring because a few labels set them badly.
The five values $(\gamma, \rho_s, C_b, C', \rho_M)$, three candidates each, are selected on pseudo-targets (Section~\ref{sec:setting-selection}).

Each step of this construction raises UAR.
Hypothesis transfer from the pooled sources with centering exceeds the best of target-only training, pooling, centered pooling, and source-weighted pooling by \NUM{p2-e31-p4-vs-best-audio-lo}--\NUM{p2-e31-p4-vs-best-audio-hi} points for audio at $K \in \{1, 5, 10\}$ (thirty iterations).
Selecting the five values adds \NUM{p2-e39-vs-p4-k5-lo}--\NUM{p2-e39-vs-p4-k5-hi} points at $K=5$ across A, A~(HuBERT), V, and AV, and \NUM{p2-e39-vs-d1-k5-lo}--\NUM{p2-e39-vs-d1-k5-hi} points over a start that weights the sources equally (thirty iterations).
We call the result the five-value start; it is the reference against which Part~II and Section~\ref{sec:comparison} measure every addition.

\subsection{Where the Direction Cannot Come From}
\label{sec:part2-limits}

Let $W^{*}$ be the student that all labels of the target pool would produce, and let $E_0 = W^{*} - W_0$ be the direction error of the start.
Two facts limit what $K$ labels can do about $E_0$.

\begin{proposition}[Span]
\label{prop:span}
Every row of the solution $W - W_0$ of \eqref{eq:student} lies in the span of the $4K$ labeled points $M^{-1}\mathbf{z}_i$, $i \in \mathcal{L}$.
\end{proposition}

The proposition is a direct consequence of the representer theorem~\cite{CITE:kimeldorf1971,CITE:scholkopf2001} applied to $W - W_0$, since the loss depends on $W$ only through the products $W\mathbf{z}_i$.
Only the projection of $E_0$ onto a subspace of at most $4K$ dimensions can be corrected, and which part of $E_0$ falls there depends on which utterances carry labels.

\begin{proposition}[Label-free design]
\label{prop:design}
Approximate the loss by a Gaussian likelihood around $W_0$ and the anchor by the prior $W \sim \mathcal{N}(W_0, C' M^{-1})$.
If $E_0$ has no preferred direction under this prior, the expected fraction of $\|E_0\|^2$ that the span of any $4K$ points captures is the same for every choice of points made without their labels.
\end{proposition}

In a Gaussian linear model the posterior covariance does not depend on the observed labels, so the best design does not depend on them either~\cite{CITE:chaloner1984,CITE:chaloner1995}.
A label-free choice can still reduce the posterior variance, for example by spreading the points, but without knowledge of the direction of $E_0$ it cannot align the corrected subspace with $E_0$.
Proposition~\ref{prop:design} rests on this approximation, so we test it directly.

We measured the fraction of $\|E_0\|^2$ captured by the span of $K=5$ labels per emotion under seven label-free selection rules (a second random order, the classifier's margin, herding on predicted classes, a k-center coreset, typicality clustering, and two designs in whitened coordinates) and under a teaching rule that knows the labels of all candidates.
Relative to random selection, the label-free rules capture \NUM{p2-e58-ratio-unlab-lo}--\NUM{p2-e58-ratio-unlab-hi} times as much, and the teaching rule \NUM{p2-e58-ratio-lll-lo}--\NUM{p2-e58-ratio-lll-hi} times as much (Fig.~\ref{fig:capture}).
Subspaces built without target labels from the disagreement of the sources or from corrections learned on pseudo-targets capture \NUM{p2-e58-r-uaub-lo}--\NUM{p2-e58-r-uaub-hi}\% of $\|E_0\|^2$.
The selection results agree: on the five-value student, no label-free rule, including standard active-learning rules, beats random selection by one point at $K=5$ (largest margin \NUM{p2-cond2-e39-maxgain-k5}), whereas replacing a weaker student by the five-value student raises random selection itself by \NUM{p2-cond2-student-random-k5-lo}--\NUM{p2-cond2-student-random-k5-hi} points.
The direction error left by a good start is thus specific to the target and invisible to label-free selection; its correction must come from label information, and from much more of it than $4K$ labels provide.

\begin{figure}[t]
\centering
\FIG{capture}{\columnwidth}
\caption{Label-free selection does not find the direction error. Fraction of the start's direction error $E_0$ captured by the span of $K=5$ labels per emotion, relative to random selection, for seven label-free rules and for a teaching rule that knows all candidate labels (upper reference). Ten iterations.}
\label{fig:capture}
\end{figure}

\subsection{Weak Labels on the Whole Pool}
\label{sec:part2-weak}

A teacher that predicts every pool utterance supplies label information on thousands of utterances.
Let $Q \in [0,1]^{N_u \times 4}$ hold the teacher's class probabilities for the $N_u$ unlabeled pool utterances.
We rescale each column of $Q$ so that every emotion receives the same total weight and scale the whole matrix to a total weight $m$; the student then minimizes \eqref{eq:student} plus $\sum_{i}\sum_{c} Q_{ic}\,\ell(W\mathbf{z}_i + b, c)$ over the pool.
The teacher's errors average out over the pool, and what moves the student is the set of utterances on which the teacher disagrees with the start.
The teacher and $m$, chosen from a small grid that includes the option of no weak labels, are selected for each $K$ on pseudo-targets.

The usefulness of a teacher has a natural scale.
Let $J$ be the UAR of the teacher minus that of the start on the pool; by definition, $J$ equals the class-averaged fraction of utterances that the teacher classifies correctly and the start incorrectly, minus the fraction for which the reverse holds.
A student that fits the teacher's labels while anchored to the start can reproduce the teacher's advantage but not create one, so the part of its gain that comes from the teacher's information should stay below $J$.
Column normalization adds a second, label-free effect: the start's own predictions, rebalanced across emotions, already act as a useful teacher.

A teacher from the same modality offers little.
Combining the votes of the per-source classifiers, each weighted by its confusion pattern on pseudo-targets, gives a belief whose $J$ over the five-value start is \NUM{p2-e52-q-minus-p-a} points for WavLM audio.
Weak labels from this belief raise the five-value start by \NUM{p2-e59k-a-k0}, \NUM{p2-e59k-a-k1}, and \NUM{p2-e59k-a-k5} points for audio at $K = 0, 1, 5$ and by \NUM{p2-e59k-v-k0}, \NUM{p2-e59k-v-k1}, and \NUM{p2-e59k-v-k5} points for the face (ten iterations), and the gain vanishes by $K=25$.
For audio at $K=0$, rebalancing the start's own predictions accounts for \NUM{p2-e59k-ctrl-a-k0} points, and the belief's information adds \NUM{p2-e59k-info-a-lo}--\NUM{p2-e59k-info-a-hi} points across $K \le 5$, below its $J$.
For text, the belief is no better than the start, and the gain is nil.

\subsection{A Teacher from Another Modality}
\label{sec:part2-crossmodal}

When the target utterances also come with other modalities, a classifier trained on the sources in those modalities knows a different direction.
This is the privileged-information setting in which the tasks share their labels and differ in their input modalities~\cite{CITE:lopezpaz2016,CITE:vapnik2009}.
For a single-modality student, the teacher is the fused start of its configuration, trained on the sources alone with the procedure of Section~\ref{sec:part2-start} and never refit on target data; the deployed student uses only its own modality.
Its $J$ is large: in the AV configuration, the fused start exceeds the audio start by \NUM{p2-e60pre-juar-av-a-fused} UAR points on the pool, and other combinations of modality predictions exceed it by \NUM{p2-e60pre-juar-av-a-lo}--\NUM{p2-e60pre-juar-av-a-hi} points, whereas the face start alone falls \NUM{p2-e60pre-juar-av-a-v} points below it.
For a fused student, combinations of single-modality predictions do not beat the fused start ($J$ between \NUM{p2-e60pre-juar-av-f-poe} and \NUM{p2-e60pre-juar-av-f-pool} points in AV), so we use a two-stage reference: the single-modality students are first trained with the fused teacher, and the product of their predictions then teaches the fused student.

Cross-modal weak labels raise every student above the five-value start (Fig.~\ref{fig:ladder}).
At $K=0$, the gains for single-modality students are \NUM{p2-e60-av-a-k0} (audio) and \NUM{p2-e60-av-v-k0} points (face) in AV, \NUM{p2-e60-vt-t-k0} (text) in VT, \NUM{p2-e60-at-t-k0} (text) in AT, and \NUM{p2-e60-avt-a-k0} (audio), \NUM{p2-e60-avt-t-k0} (text), and \NUM{p2-e60-avt-v-k0} points (face) in AVT; all \NUM{p2-e60-m1-n} single-modality students gain at both $K=0$ and $K=5$.
The fused students gain \NUM{p2-e60-f-k0-lo}--\NUM{p2-e60-f-k0-hi} points at $K=0$.
These gains exceed the gain from a teacher of the same modality and remain below the $J$ of the cross-modal teacher, as the scale of the previous subsection anticipates.
With more labels the gains shrink, and at $K \ge 10$ \NUM{p2-e60-m3-loss-n} of the \NUM{p2-e60-m3-students-n} students fall slightly below the five-value start (at most \NUM{p2-e60-m3-worst} points).
The extended family, which adds as candidate teachers the fused student refit on the $K$ labels and the two-stage reference and adds the weighting of the weak labels (balanced, unbalanced, or hard) to the selection, removes these losses (Section~\ref{sec:comparison}).

\begin{figure}[t]
\centering
\FIG{ladder}{\columnwidth}
\caption{The method step by step. UAR of target-only training, pooling, hypothesis transfer from the pooled sources, the five-value start, weak labels from a same-modality belief, and weak labels from the fused teacher, against $K$, on a common set of targets.}
\label{fig:ladder}
\end{figure}

\section{Comparison with Standard Methods}
\label{sec:comparison}


\subsection{Competitors and Protocol}
\label{sec:comparison-protocol}

The nine competitors are self-training on the student's own confident predictions in the form of Lee~\cite{lee2013pseudolabel}, hard pseudo-labels from the teacher above a confidence threshold, entropy minimization~\cite{grandvalet2004semisupervised}, distillation with a temperature~\cite{hinton2015distilling}, generalized distillation~\cite{lopezpaz2016unifying}, co-training~\cite{blum1998combining}, co-regularization~\cite{sindhwani2005coregularization}, adjustment of the class priors by expectation maximization~\cite{saerens2002adjusting}, and rebalancing of the start's outputs.
Every method that retrains the classifier trains the same anchored linear student \eqref{eq:student} from the same five-value start on the same $K$ labels and unlabeled pool; prior adjustment and rebalancing correct the outputs of the start. Competitors that need a teacher receive the fused start (distillation with a temperature and hard pseudo-labels for single-modality students), the fused student refit on the $K$ labels (generalized distillation), or the product of the single-modality starts (distillation for fused students); none receives a two-stage reference, so the margins include the choice of teacher. For fused students, generalized distillation and hard pseudo-labels are not defined, and co-training enters in two variants.
The \emph{tuned} version of each competitor selects its hyperparameters (thresholds, temperatures, weights of the unlabeled loss) on the same pseudo-targets and for each $K$, as our method does; the \emph{fixed} version uses the values of the original papers.
Ours denotes the basic method of Section~\ref{sec:part2-crossmodal} and ours-ext the extended family.
The AVT configuration evaluates the audio, face, text, and fused students on \NUM{cmp-targets-avt-n} targets, the AV, VT, and AT configurations the fused student on \NUM{cmp-targets-av-n}, \NUM{cmp-targets-vt-n}, and \NUM{cmp-targets-at-n} targets, and a second AVT configuration replaces every encoder by its second encoder (HuBERT, CLIP, DeBERTa); results average \NUM{cmp-reps-n} iterations per target, and pseudo-targets use \NUM{cmp-pseudo-reps-n}.

\subsection{Results}
\label{sec:comparison-results}

Ours-ext exceeds the strongest tuned competitor for every student at $K=0$ and $K=5$, by \NUM{cmp-oursx-vs-best-range-lo}--\NUM{cmp-oursx-vs-best-range-hi} points (Table~\ref{tab:compare}; Fig.~S6 of the supplementary material shows the range across targets).
No single competitor is strongest everywhere.
The basic method, whose candidate teachers do not include the fused student refit on the $K$ labels, exceeds the strongest competitor in \NUM{cmp-ours-vs-best-pos-n} of the \NUM{cmp-ours-vs-best-cells-n} combinations of student and $K \in \{0, 5\}$; at $K=5$ it falls behind generalized distillation for the AVT audio and text students and is level with co-training for the AT fused student.
Across label budgets (Fig.~S7), ours-ext keeps its advantage for the audio student at every $K$ (\NUM{cmp-oursx-vs-best-avt-a-lo}--\NUM{cmp-oursx-vs-best-avt-a-hi} points), and generalized distillation overtakes it only for the text student at $K \ge 10$, by \NUM{cmp-gd-vs-oursx-avt-t-kge10-lo}--\NUM{cmp-gd-vs-oursx-avt-t-kge10-hi} points.
Over the five-value start, ours-ext gains \NUM{cmp-oursx-vs-e39-avt-a-lo}--\NUM{cmp-oursx-vs-e39-avt-a-hi} points for the AVT audio student at every $K$, \NUM{cmp-oursx-vs-e39-avt-t-kle5-lo}--\NUM{cmp-oursx-vs-e39-avt-t-kle5-hi} for text at $K \le 5$, and \NUM{cmp-oursx-vs-e39-avt-f-kge1-lo}--\NUM{cmp-oursx-vs-e39-avt-f-kge1-hi} for the fused student at $K \ge 1$, without the losses of the basic method at $K \ge 10$.
We test each margin over target and iteration pairs with a two-sided Wilcoxon signed-rank test, with the Holm correction over all students and budgets.
The margin over the strongest competitor is positive in \NUM{st-a-k0-pos-n} of \NUM{st-a-k0-cells-n} students at $K=0$ and in \NUM{st-a-k5-pos-n} of \NUM{st-a-k5-cells-n} at $K=5$, with effect sizes $d_z$ (mean difference over its standard deviation) of \NUM{st-a-k0-dz-lo}--\NUM{st-a-k0-dz-hi} and \NUM{st-a-k5-dz-lo}--\NUM{st-a-k5-dz-hi}, and it is significant at the $5\%$ level for \NUM{st-a-k0-sig-n} and \NUM{st-a-k5-sig-n} students.
Because several iterations of a target share a test fold, we also average the iterations that share a fold; this leaves \NUM{st-a-k0-sigdedup-n} and \NUM{st-a-k5-sigdedup-n} significant students (at $K=0$ the AVT audio student and the VT and AT fused students; at $K=5$ all but the AVT face and text students and the AV and AT fused students).
With four to six targets per configuration, these tests describe the variation within our corpora, not generalization to other corpora.

\begin{table*}[t]
\caption{Comparison under one protocol: UAR of the five-value start and differences from it (points) for the strongest tuned competitor, ours, and ours-ext, averaged over targets and ten iterations. ``None'' means that every tuned competitor selected no use of the unlabeled data. Competitors: GD, generalized distillation; KD, distillation with a temperature; HPL, hard pseudo-labels; RB, output rebalancing; ST, self-training; CoT, co-training; CoR, co-regularization.}
\label{tab:compare}
\centering
\footnotesize
\setlength{\tabcolsep}{4pt}
\begin{tabular}{@{}llrrlrrrrrr@{}}
\toprule
 & & \multicolumn{2}{c}{Start (UAR)} & Strongest & \multicolumn{2}{c}{Strongest $-$ start} & \multicolumn{2}{c}{Ours $-$ start} & \multicolumn{2}{c}{Ours-ext $-$ start} \\
\cmidrule(lr){3-4}\cmidrule(lr){6-7}\cmidrule(lr){8-9}\cmidrule(l){10-11}
Config. & Student & $K{=}0$ & $K{=}5$ & ($K{=}0$/$5$) & $K{=}0$ & $K{=}5$ & $K{=}0$ & $K{=}5$ & $K{=}0$ & $K{=}5$ \\
\midrule
AVT & A & \NUM{tab-avt-a-k0-e39} & \NUM{tab-avt-a-k5-e39} & GD/GD & \NUM{cmp-best-vs-e39-avt-a-k0} & \NUM{cmp-best-vs-e39-avt-a-k5} & \NUM{cmp-ours-vs-e39-avt-a-k0} & \NUM{cmp-ours-vs-e39-avt-a-k5} & \NUM{cmp-oursx-vs-e39-avt-a-k0} & \NUM{cmp-oursx-vs-e39-avt-a-k5} \\
 & V & \NUM{tab-avt-v-k0-e39} & \NUM{tab-avt-v-k5-e39} & RB/none & \NUM{cmp-best-vs-e39-avt-v-k0} & \NUM{cmp-best-vs-e39-avt-v-k5} & \NUM{cmp-ours-vs-e39-avt-v-k0} & \NUM{cmp-ours-vs-e39-avt-v-k5} & \NUM{cmp-oursx-vs-e39-avt-v-k0} & \NUM{cmp-oursx-vs-e39-avt-v-k5} \\
 & T & \NUM{tab-avt-t-k0-e39} & \NUM{tab-avt-t-k5-e39} & HPL/GD & \NUM{cmp-best-vs-e39-avt-t-k0} & \NUM{cmp-best-vs-e39-avt-t-k5} & \NUM{cmp-ours-vs-e39-avt-t-k0} & \NUM{cmp-ours-vs-e39-avt-t-k5} & \NUM{cmp-oursx-vs-e39-avt-t-k0} & \NUM{cmp-oursx-vs-e39-avt-t-k5} \\
 & Fused & \NUM{tab-avt-f-k0-e39} & \NUM{tab-avt-f-k5-e39} & RB/ST & \NUM{cmp-best-vs-e39-avt-f-k0} & \NUM{cmp-best-vs-e39-avt-f-k5} & \NUM{cmp-ours-vs-e39-avt-f-k0} & \NUM{cmp-ours-vs-e39-avt-f-k5} & \NUM{cmp-oursx-vs-e39-avt-f-k0} & \NUM{cmp-oursx-vs-e39-avt-f-k5} \\
AV & Fused & \NUM{tab-av-f-k0-e39} & \NUM{tab-av-f-k5-e39} & RB/ST & \NUM{cmp-best-vs-e39-av-f-k0} & \NUM{cmp-best-vs-e39-av-f-k5} & \NUM{cmp-ours-vs-e39-av-f-k0} & \NUM{cmp-ours-vs-e39-av-f-k5} & \NUM{cmp-oursx-vs-e39-av-f-k0} & \NUM{cmp-oursx-vs-e39-av-f-k5} \\
VT & Fused & \NUM{tab-vt-f-k0-e39} & \NUM{tab-vt-f-k5-e39} & RB/ST & \NUM{cmp-best-vs-e39-vt-f-k0} & \NUM{cmp-best-vs-e39-vt-f-k5} & \NUM{cmp-ours-vs-e39-vt-f-k0} & \NUM{cmp-ours-vs-e39-vt-f-k5} & \NUM{cmp-oursx-vs-e39-vt-f-k0} & \NUM{cmp-oursx-vs-e39-vt-f-k5} \\
AT & Fused & \NUM{tab-at-f-k0-e39} & \NUM{tab-at-f-k5-e39} & RB/CoT & \NUM{cmp-best-vs-e39-at-f-k0} & \NUM{cmp-best-vs-e39-at-f-k5} & \NUM{cmp-ours-vs-e39-at-f-k0} & \NUM{cmp-ours-vs-e39-at-f-k5} & \NUM{cmp-oursx-vs-e39-at-f-k0} & \NUM{cmp-oursx-vs-e39-at-f-k5} \\
\midrule
AVT, & A & \NUM{tab-w2-a-k0-e39} & \NUM{tab-w2-a-k5-e39} & CoT/GD & \NUM{cmp-best-vs-e39-w2-a-k0} & \NUM{cmp-best-vs-e39-w2-a-k5} & \NUM{cmp-ours-vs-e39-w2-a-k0} & \NUM{cmp-ours-vs-e39-w2-a-k5} & \NUM{cmp-oursx-vs-e39-w2-a-k0} & \NUM{cmp-oursx-vs-e39-w2-a-k5} \\
second & V & \NUM{tab-w2-v-k0-e39} & \NUM{tab-w2-v-k5-e39} & CoT/KD & \NUM{cmp-best-vs-e39-w2-v-k0} & \NUM{cmp-best-vs-e39-w2-v-k5} & \NUM{cmp-ours-vs-e39-w2-v-k0} & \NUM{cmp-ours-vs-e39-w2-v-k5} & \NUM{cmp-oursx-vs-e39-w2-v-k0} & \NUM{cmp-oursx-vs-e39-w2-v-k5} \\
encoders & T & \NUM{tab-w2-t-k0-e39} & \NUM{tab-w2-t-k5-e39} & RB/GD & \NUM{cmp-best-vs-e39-w2-t-k0} & \NUM{cmp-best-vs-e39-w2-t-k5} & \NUM{cmp-ours-vs-e39-w2-t-k0} & \NUM{cmp-ours-vs-e39-w2-t-k5} & \NUM{cmp-oursx-vs-e39-w2-t-k0} & \NUM{cmp-oursx-vs-e39-w2-t-k5} \\
 & Fused & \NUM{tab-w2-f-k0-e39} & \NUM{tab-w2-f-k5-e39} & RB/CoT & \NUM{cmp-best-vs-e39-w2-f-k0} & \NUM{cmp-best-vs-e39-w2-f-k5} & \NUM{cmp-ours-vs-e39-w2-f-k0} & \NUM{cmp-ours-vs-e39-w2-f-k5} & \NUM{cmp-oursx-vs-e39-w2-f-k0} & \NUM{cmp-oursx-vs-e39-w2-f-k5} \\
\bottomrule
\end{tabular}
\end{table*}



\subsection{Fixed Hyperparameters and Entropy Minimization}
\label{sec:comparison-fixed}

With the hyperparameters of the original papers, entropy minimization, co-regularization, and prior adjustment change the UAR of the single-modality students of AVT at $K=0$ by \NUM{cmp-fixed-collapse-avt-mod-k0-lo} to \NUM{cmp-fixed-collapse-avt-mod-k0-hi} points, and co-regularization, and entropy minimization for the audio student, reduce it to chance level.
For the fused students, fixed entropy minimization changes it by \NUM{cmp-entmin-fixed-vs-e39-f-k0-lo} to \NUM{cmp-entmin-fixed-vs-e39-f-k0-hi} points on average over the four configurations, depending on the target (Fig.~S8): by $+$\NUM{cmp-entmin-fixed-vs-e39-avt-f-k0-impro} (AVT) and $+$\NUM{cmp-entmin-fixed-vs-e39-at-f-k0-impro} (AT) on the improvised IEMOCAP sessions, but by \NUM{cmp-entmin-fixed-vs-e39-vt-f-k0-impro} (VT) and \NUM{cmp-entmin-fixed-vs-e39-av-f-k0-impro} (AV) on the same target and by \NUM{cmp-entmin-fixed-vs-e39-avt-f-k0-meld} on MELD in AVT.
Pseudo-target selection detects these failures and switches entropy minimization off for the AVT fused student at every $K$.


\subsection{Second Encoders, Fine-Tuning, and Selection}
\label{sec:comparison-robust}

With the second encoders, the basic method raises every student above the start at $K=0$ and $K=5$ (Table~\ref{tab:compare}), and ours-ext exceeds the strongest tuned competitor by \NUM{cmp-oursx-vs-best-w2-k0k5-lo}--\NUM{cmp-oursx-vs-best-w2-k0k5-hi} points at $K=0$ and $K=5$ and by \NUM{cmp-oursx-vs-best-w2-k25-lo}--\NUM{cmp-oursx-vs-best-w2-k25-hi} points at $K=25$.

To test whether the frozen setting understates what adaptation can do, we fine-tuned WavLM with low-rank adapters on the audio student of four targets over three iterations (Fig.~S9 and Section~S7 of the supplementary material).
Fine-tuning on $K=5$ labels per emotion falls \NUM{cmp-ft-konly-below-frozen-e39-k5} points below the frozen five-value start (\NUM{cmp-ft-konly-below-frozen-e39-k5-tgt-lo}--\NUM{cmp-ft-konly-below-frozen-e39-k5-tgt-hi} across targets).
The weak labels of ours-ext, fixed to the form it selects for the AVT audio student, still help the fine-tuned model, by \NUM{cmp-ft-weak-vs-start-k0} points at $K=0$ and \NUM{cmp-ft-weak-vs-anchor-k5} at $K=5$, but the fine-tuned version of our method stays \NUM{cmp-ft-ours-below-frozen-ours-k0} and \NUM{cmp-ft-ours-below-frozen-ours-k5} points below the same configuration on frozen features, and centering closes only \NUM{cmp-ftc-gap-closed-start-k0}--\NUM{cmp-ftc-gap-closed-konly-k5}\% of the gap.


\emph{Loss due to selection.}
\label{sec:comparison-selection}
Pseudo-target selection does not always find the best candidate of our family: at $K=0$, the best candidate on the test speakers exceeds the selected one by \NUM{cmp-oursx-oracle-vs-tuned-avt-t-k0} points for the AVT text student and by \NUM{cmp-oursx-oracle-vs-tuned-w2-t-k0} points for the text student with second encoders, and by \NUM{cmp-oursx-oracle-vs-tuned-at-f-k0}--\NUM{cmp-oursx-oracle-vs-tuned-w2-v-k0} points for the other students.
The gap reaches one point for \NUM{cmp-oursx-oracle-gap-ge1-n} of the \NUM{cmp-oursx-oracle-gap-students-n} students; it bounds what a better selection rule could add, and because the competitors are selected by the same procedure, it does not favor our method.

\section{Discussion}
\label{sec:discussion}

\emph{Practical guidance.}
The results translate into four recommendations for a new corpus with a few labels per emotion.
First, center the source and the target features separately; this uses no label and removes most of the threshold loss that a source-only classifier carries (Section~\ref{sec:part1-fix}).
Second, start from all sources and spend the $K$ labels on the intercepts with a loose anchor while anchoring the weights strongly, since the labels set the intercepts well and the weights badly (Section~\ref{sec:part2-start}).
Third, do not expect a label-free choice of which utterances to label to repair the direction of a good start (Section~\ref{sec:part2-limits}).
Fourth, if other modalities are recorded with the target utterances, even only at labeling time, let a fused teacher trained on the sources with those modalities label the whole pool; the deployed classifier can remain single-modality (Sections~\ref{sec:part2-crossmodal} and~\ref{sec:comparison}).

\emph{Why text behaves differently.}
Text features lose little through the threshold and much through the direction (Table~\ref{tab:part1}), so the few labels that fix thresholds elsewhere have little threshold to fix in text, and for text alone the threshold share of the gain is at most \NUM{p1-e29-critshare-k1-t}.
The same two-cost account covers this case: the share of the gain that comes from the threshold follows the share of the loss that the threshold causes.
For the method, text students gain from the fused teacher as much as audio students do, because the fused teacher supplies the direction information that text labels lack.

\emph{Entropy minimization and the cluster assumption.}
Entropy minimization pushes decision boundaries away from dense regions of unlabeled data and helps when the classes form separate clusters.
In our data it helps on the improvised IEMOCAP sessions in the AVT and AT configurations and fails badly on MELD and CMU-MOSEI and on the same IEMOCAP target in the VT and AV configurations (Section~\ref{sec:comparison-fixed}).
The cross-modal weak labels rely on a different property, the complementarity between modalities, and their gains vary much less: across the \NUM{st-b-tgt-n} combinations of target and cell at $K=0$ and $K=5$, ours-ext falls below the start in \NUM{st-b-tgtneg-n}, with a worst difference of \NUM{st-b-tgtneg-worst} points.

\emph{What a source is.}
Part~I gives the value of a source in labels of direction, and Part~II shows that, in the cases we measured, the gain from a teacher stays below its UAR advantage over the start.
Here $q$, $e$, and $J$ were measured with target labels; whether pseudo-targets estimate them well enough to choose, before deployment, between more labels, another source, and another modality remains to be tested.

\section{Limitations and Ethics}
\label{sec:limitations-ethics}
% TODO: about 300 words; frozen features, linear head, four emotions, seven corpora; extrapolation limits; ethics of secondary use of public corpora.

\section{Conclusion}
\label{sec:conclusion}

With a few target labels per emotion and a linear head on frozen features, the labels first fix the decision thresholds: for single-modality audio and face at $K \le 5$, \NUM{p1-e29-critshare-kle5-single-lo}--\NUM{p1-e29-critshare-kle5-single-hi} of the improvement that pooling brings comes from the threshold, while the direction of audio classifiers improves only at the rate $q/K$ with $q$ between \NUM{p1-qdir-audio-lo} and \NUM{p1-qdir-audio-hi}.
A source supplies \NUM{p1-e32-n1src-audio-lo}--\NUM{p1-e32-n1src-audio-hi} labels per emotion of direction for audio, and neither a few labels nor a label-free choice of the utterances to label moves the direction of a good start much.
Weak labels on the whole unlabeled pool from a teacher of another modality, trained on the sources alone, move it: they raise every single-modality student above the five-value start, and the extended method exceeds the strongest competitor tuned under the same protocol in every configuration and with second encoders.
The two costs and the teacher's UAR advantage over the start indicate in advance where a few labels, a source, or another modality pays off.


\bibliographystyle{IEEEtran}
\bibliography{references}

\end{document}
